\documentclass[a4paper,10pt]{article}


% 载入multicol包用于创建多列布局
\usepackage{multicol}
% 载入amsmath包用于数学公式
\usepackage{amsmath}
\usepackage{amssymb}
% 设置页边距
\usepackage{soul}
\usepackage{xcolor}
\sethlcolor{cyan}
\usepackage[margin=0.5in]{geometry}
% 载入graphicx包用于插入图片
\usepackage{graphicx}
% 设置字体大小
\usepackage{times}
\renewcommand{\baselinestretch}{0.9} % 行距

% 开始文档
\begin{document}

% 开启两列布局，调整全局字体大小为小一点
\begin{multicols}{2}
\small

% 第一列开始

% 调整section标题字号
\section*{Set Theory}
\begin{itemize}
    \item Def: Set
    \item Def: Subset, proper subset
    \item Def: Union, intersection, complement, difference, disjoint
    \item Theorem (DeMorgan):\\ $(B \cup C)^c = B^c \cap C^c, \quad (B \cap C)^c = B^c \cup C^c$\\
    $A \setminus (B \cup C) = (A \setminus B) \cap (A \setminus C), \quad 
    A \setminus (B \cap C) = (A \setminus B) \cup (A \setminus C).$
    \item well ordering property of N: Every nonempty subset of N has a least element.
    \item Theorem (Principle of induction)
    \item Def: Cartesian product
    \item Def: Function (domain $D(f)$, range $R(f)$, codomain)
    \item Def: Image $f(C)$, inverse image $f^{-1}(D)$
    \item \hl{Def: Injection, Surjection, Bijection}
    \item Def: composition of functions
    \item Def: Binary relation (eg. <, =)
    \item \hl{Def: Equivalence relation: Reflexive, Symmetric, Transitive}
    \item Equivalence relation eg.[a], [(a,b)] or a/b\\
    eg. For equivalence relation $\sim$ on N, $[n]=\{m:m\sim n\},N/\sim=\{[n]\}$
    \item Def: Cardinality $|A|=|B|\Leftrightarrow$ exist a bijection $f:A\to B$
    \item Def: $|A|=n$: A has same cardinality as {1,2,...,n}
    \item Def: $|A|\le |B|$: exist an injection from A to B
    \item \hl{Def: Countably (infinite): $|A|=|N|$}
    \item Def: Power set
    \item \hl{Theorem (Cantor): $|A|<|P(A)|$}
    
\end{itemize}

\section*{Real Numbers}
\begin{itemize}
    \item \hl{Def: Ordered set:} trichotomy, transitivity
    \item \hl{Def: upper/lower bound, least upper bound (supremum), greatest lower bound (infimum)}
    \item \hl{Def: least-upper-bound property:} For nonempty $E\subset S$ bounded above, sup E exsits in S. (R, not Q)
    \item \hl{Def: Field:} \\
    $1. x+y\in F,\\ 
    2. x+y=y+x;\\ 
    3. (x+y)+z=x+(y+z)\\ 
    4. \exists 0+x=x\\
    5. \exists -x+x=0\\ 
    6. xy\in F\\
    7. xy=yx\\
    8. (xy)z=x(yz)\\
    9. \exists 1x=x\\
    10. \exists x(1/x)=1\\
    11. x(y+z)=xy+xz$
    \item \hl{Def: Ordered field:} F is an ordered set s.t.\\
    (1) $x<y\Rightarrow x+z<y+z$\\
    (2) $x>0, y>0 \Rightarrow xy>0$
    \item Prop: $x<\forall \epsilon>0\Rightarrow x\le 0$
    \item Theorem (Archimedean Property) $\exists n \in N,s.t.nx>y$
    \item Theorem (Q dense in R): $\exists r\in Q, s.t. x<r<y$
    \item Def: $R*=R\cup \{-\infty,\infty\}$
    \item Def: Maxima, Minima
    \item Def: Absolute value
    \item Prop: (Triangle inequality): $|x+y|\le|x|+|y|$\\
    $||x|-|y||\le |x-y|$, $|x-y|\le |x|+|y|$
    \item \hl{Def: function bounded:} $|f(x)|\le M$
    \item Def: Interval, closed/open interval
    \item Theorem (Cantor): R is uncountable.
    
\end{itemize}

\section*{Sequences and Series}
\begin{itemize}
    \item Def: Sequence, bounded sequence
    \item \hl{Def: Sequence convergent/divergent}
    \item Prop: Convergent sequence is bounded.
    \item Def: monotone increasing/decreasing sequence
    \item \hl{Theorem (Monotone converge theorem):} monotone sequence bounded $\Leftrightarrow$ convergent
    \item Def: K-tail of sequence
    \item Def: subsequence
    \item \hl{Prop: sequence convergent $\Rightarrow$ all subsequences converge to same value.}
    \item Lemma: Squeeze lemma
    \item Lemma: $x_n\le y_n\Rightarrow lim_{n\to \infty}x_n\le lim_{n\to \infty}y_n$
    \item Algebraic operations of limits ($+,-,*,/,sqrt,||$)
    \item Lemma (Ratio test for sequences): If $L=lim_{n\to \infty} \frac{x_{n+1}}{x_{n}}$ exists, $L<1\Rightarrow lim_{n\to \infty}x_n=0$; $L>1\Rightarrow x_n$ unbounded.
    \item \hl{Def: $x_n$ bounded:}\\ $lim_{n\to \infty}sup\ x_n=lim_{n\to \infty}\ sup\{x_k:k\ge n\}=lim_{n\to \infty}a_n$,\\ $lim_{n\to \infty}inf\ x_n=lim_{n\to \infty}\ inf\{x_k:k\ge n\}=lim_{n\to \infty}b_n$
    \item Theorem: ${x_n}$ bounded:\\ $\lim_{k \to \infty} x_{n_k} = \limsup_{n \to \infty} x_n$,\\ $\lim_{k \to \infty} x_{m_k} = \liminf_{n \to \infty} x_n$
    \item \hl{Prop}: $x_n$ bounded: $x_n\ converges \Leftrightarrow \limsup_{n \to \infty} x_n=\liminf_{n \to \infty} x_n=lim_{n\to \infty}x_n$
    \item Prop: $x_n$ bounded: $x_n$ converges to x $\Leftrightarrow$ every convergent subsequence converges to x.
    \item \hl{Theorem (Bolzano-Weierstrass)}: $x_n$ bounded: exist convergent subsequence 
    $\{x_{n_i}\}_{i=1}^{\infty}$
    \item Def: Diverge to infinity/minus infinity
    \item Prop: Monotone unbounded sequence diverges to infinity/minus infinity.
    \item \hl{Def: Cauchy sequence}
    \item Prop: If a sequence is Cauchy, then it is bounded.
    \item \hl{Theorem: A sequence of real numbers is Cauchy $\Leftrightarrow$ it converges.}
    \item Def: series $\sum_{n=1}^{\infty}x_n=lim_{k\to \infty}s_k$
    \item Def: $\sum_{n=1}^{\infty}x_n$ is Cauchy if the partial sums $\{s_n\}_{n=1}^\infty$ is a Cauchy sequence.
    \item Prop: $\sum_{n=1}^{\infty} x_n$ \textit{is Cauchy if and only if for every} $\epsilon > 0$, \textit{there exists an} $M \in \mathbb{N}$ \textit{such that for every} $n \geq M$ \textit{and every} $k > n$, $\left| \sum_{i=n+1}^{k} x_i \right| < \epsilon.$
    \item Prop: \textit{Let} $\sum_{n=1}^{\infty} x_n$ \textit{be a convergent series. Then the sequence} $\{x_n\}_{n=1}^{\infty}$ \textit{is convergent and} $\lim_{n \to \infty} x_n = 0.$
    \item Prop: Linearity of series
    \item Prop: \textit{If} $x_n \geq 0$ \textit{for all} $n$, \textit{then} $\sum_{n=1}^{\infty} x_n$ \textit{converges if and only if the sequence of partial sums is bounded above.}
    \item Def: Converge absolutly, converge conditionally
    \item Prop: $\sum_{n=1}^{\infty} x_n$ \textit{converges absolutely, then it converges.}
    \item Prop: Comparison test
    \item Prop: p-series/ p-test: $\sum_{n=1}^{\infty} \frac{1}{n^p}$ converges if and only if $p>1$.
    \item Prop: (Ratio test). \textit{Let} $\sum_{n=1}^{\infty} x_n$ \textit{be a series,} $x_n \neq 0$ , $L := \lim_{n \to \infty} \frac{|x_{n+1}|}{|x_n|} \quad \textit{exists}.$  \textit{If} $L < 1$, \textit{then} $\sum_{n=1}^{\infty} x_n$ \textit{converges absolutely}. \textit{If} $L > 1$, \textit{then} $\sum_{n=1}^{\infty} x_n$ \textit{diverges}.
    \item Prop (Root test): Let $\sum_{n=1}^{\infty} x_n$ be a series and let $L := \limsup_{n \to \infty} |x_n|^{1/n}.$ (i) If $L < 1$, then $\sum_{n=1}^{\infty} x_n$ converges absolutely.  
    (ii) If $L > 1$, then $\sum_{n=1}^{\infty} x_n$ diverges.
    \item Prop (Alternating series): Let $\{x_n\}_{n=1}^{\infty}$ be a monotone decreasing sequence of positive real numbers such that $\lim\limits_{n \to \infty} x_n = 0$. Then $\sum_{n=1}^{\infty} (-1)^n x_n$ converges.
    \item Prop: Let $\sum_{n=1}^{\infty} x_n$ be an absolutely convergent series converging to a number $x$. Let $\sigma: \mathbb{N} \to \mathbb{N}$ be a bijection. Then $\sum_{n=1}^{\infty} x_{\sigma(n)}$ is absolutely convergent and converges to $x$.
    \item Theorem (Mertens’ theorem). Suppose $\sum_{n=0}^{\infty} a_n$ and $\sum_{n=0}^{\infty} b_n$ are two convergent series, converging to $A$ and $B$ respectively. Suppose at least one of the series converges absolutely. Define $c_n := a_0 b_n + a_1 b_{n-1} + \cdots + a_n b_0 = \sum_{i=0}^{n} a_i b_{n-i}.$ Then the series $\sum_{n=0}^{\infty} c_n$ converges to $AB$. The series $\sum_{n=0}^{\infty} c_n$ is called the \textit{Cauchy product}.
    \item Prop: Let $\sum_{n=0}^{\infty} a_n (x - x_0)^n$ be a power series. If the series is convergent, then either it converges absolutely at all $x \in \mathbb{R}$, or there exists a number $\rho$, such that the series converges absolutely on the interval $(x_0 - \rho, x_0 + \rho)$ and diverges when $x < x_0 - \rho$ or $x > x_0 + \rho$.
    \item Prop: Let $\sum_{n=0}^{\infty} a_n (x - x_0)^n$ be a power series, and let $R := \limsup_{n \to \infty} |a_n|^{1/n}.$ If $R = \infty$, the power series is divergent. If $R = 0$, then the power series converges everywhere. Otherwise, the radius of convergence $\rho = 1/R$.
    

\end{itemize}

\section*{Continuous Functions}
\begin{itemize}
    \item \hl{Def (cluster point)}: Let \( S \subset \mathbb{R} \) be a set. A number \( x \in \mathbb{R} \) is called a \textit{cluster point} of \( S \) if for every \( \epsilon > 0 \), the set \( (x - \epsilon, x + \epsilon) \cap S \setminus \{x\} \) is not empty.
    \item Prop: \textit{Let \( S \subset \mathbb{R} \). Then \( x \in \mathbb{R} \) is a cluster point of \( S \) if and only if there exists a convergent sequence of numbers \( \{x_n\}_{n=1}^{\infty} \) such that \( x_n \neq x \) and \( x_n \in S \) for all \( n \), and} $\lim_{n \to \infty} x_n = x.$
    \item \hl{Def (limit of function)}: Let \( f: S \to \mathbb{R} \) be a function and \( c \) a cluster point of \( S \subset \mathbb{R} \). Suppose there exists an \( L \in \mathbb{R} \) and for every \( \epsilon > 0 \), there exists a \( \delta > 0 \) such that whenever \( x \in S \setminus \{c\} \) and \( |x - c| < \delta \), we have\[|f(x) - L| < \epsilon.\]We then say \( f(x) \) \textit{converges} to \( L \) as \( x \) goes to \( c \), and we write\[f(x) \to L \quad \text{as} \quad x \to c.\]We say \( L \) is a \textit{limit} of \( f(x) \) as \( x \) goes to \( c \), and if \( L \) is unique (it is), we write\[\lim_{x \to c} f(x) := L.\]
    \item Prop: \textit{Let \( c \) be a cluster point of \( S \subset \mathbb{R} \) and let \( f: S \to \mathbb{R} \) be a function such that \( f(x) \) converges as \( x \) goes to \( c \). Then the limit of \( f(x) \) as \( x \) goes to \( c \) is unique.}
    \item Lemma: \textit{Let \( S \subset \mathbb{R} \), let \( c \) be a cluster point of \( S \), let \( f: S \to \mathbb{R} \) be a function, and let \( L \in \mathbb{R} \). Then \( f(x) \to L \) as \( x \to c \) if and only if for every sequence \( \{x_n\}_{n=1}^{\infty} \) such that \( x_n \in S \setminus \{c\} \) for all \( n \), and such that \( \lim_{n \to \infty} x_n = c \), we have that the sequence \( \{f(x_n)\}_{n=1}^{\infty} \) converges to \( L \).}
    \item Corollary: $f(x) \leq g(x) \quad \text{for all } x \in S \setminus \{c\}\Rightarrow \lim_{x \to c} f(x) \leq \lim_{x \to c} g(x).$
    \item Corollary (Squeeze Thm): \textit{Let \( S \subset \mathbb{R} \) and let \( c \) be a cluster point of \( S \). Suppose \( f: S \to \mathbb{R} \), \( g: S \to \mathbb{R} \), and \( h: S \to \mathbb{R} \) are functions such that}
    \[
    f(x) \leq g(x) \leq h(x) \quad \text{for all } x \in S \setminus \{c\}.
    \]
    \textit{Suppose the limits of \( f(x) \) and \( h(x) \) as \( x \) goes to \( c \) both exist, and}
    \[
    \lim_{x \to c} f(x) = \lim_{x \to c} h(x).
    \]
    \textit{Then the limit of \( g(x) \) as \( x \) goes to \( c \) exists and}
    \[
    \lim_{x \to c} g(x) = \lim_{x \to c} f(x) = \lim_{x \to c} h(x).
    \]
    \item Corollary: Algebraic operations of function limits $(+,-,*,/,||)$
    \item \hl{Def (continuous function)}: Suppose \( S \subset \mathbb{R} \) and \( c \in S \). We say \( f: S \to \mathbb{R} \) is \textit{continuous at} \( c \) if for every \( \epsilon > 0 \) there is a \( \delta > 0 \) such that whenever \( x \in S \) and \( |x - c| < \delta \), we have \( |f(x) - f(c)| < \epsilon \). 

    \item Prop: When \( f: S \to \mathbb{R} \) is continuous at all \( c \in S \), then we simply say \( f \) is a \textit{continuous function}.
    \begin{itemize}
        \item[(i)] \textit{If \( c \) is not a cluster point of \( S \), then \( f \) is continuous at \( c \).}
        
        \item[(ii)] \textit{If \( c \) is a cluster point of \( S \), then \( f \) is continuous at \( c \) if and only if the limit of \( f(x) \) as \( x \to c \) exists and}
        \[
        \lim_{x \to c} f(x) = f(c).
        \]
        
        \item[(iii)] \textit{The function \( f \) is continuous at \( c \) if and only if for every sequence \( \{x_n\}_{n=1}^{\infty} \) where \( x_n \in S \) and \( \lim\limits_{n \to \infty} x_n = c \), the sequence \( \{ f(x_n) \}_{n=1}^{\infty} \) converges to \( f(c) \).}
    \end{itemize}
    \item Prop: \textit{Let \( f: \mathbb{R} \to \mathbb{R} \) be a polynomial. That is,}
    \[
    f(x) = a_d x^d + a_{d-1} x^{d-1} + \cdots + a_1 x + a_0,
    \]
    \textit{for some constants \( a_0, a_1, \dots, a_d \). Then \( f \) is continuous.}

    \item Prop: \textit{Let \( f: S \to \mathbb{R} \) and \( g: S \to \mathbb{R} \) be functions continuous at \( c \in S \).}
    \begin{itemize}
        \item[(i)] \textit{The function \( h: S \to \mathbb{R} \) defined by \( h(x) := f(x) + g(x) \) is continuous at \( c \).}
        
        \item[(ii)] \textit{The function \( h: S \to \mathbb{R} \) defined by \( h(x) := f(x) - g(x) \) is continuous at \( c \).}
        
        \item[(iii)] \textit{The function \( h: S \to \mathbb{R} \) defined by \( h(x) := f(x) g(x) \) is continuous at \( c \).}
        
        \item[(iv)] \textit{If \( g(x) \neq 0 \) for all \( x \in S \), the function \( h: S \to \mathbb{R} \) given by \( h(x) := \frac{f(x)}{g(x)} \) is continuous at \( c \).}
    \end{itemize}
    \item Prop: \textit{Let \( A, B \subset \mathbb{R} \) and \( f: B \to \mathbb{R} \) and \( g: A \to B \) be functions. If \( g \) is continuous at \( c \in A \) and \( f \) is continuous at \( g(c) \), then \( f \circ g: A \to \mathbb{R} \) is continuous at \( c \).}
    \item \hl{Prop}: \textit{Let \( f: S \to \mathbb{R} \) be a function and \( c \in S \). Suppose there exists a sequence \( \{ x_n \}_{n=1}^{\infty} \), \( x_n \in S \) for all \( n \), and \( \lim\limits_{n \to \infty} x_n = c \) such that \( \{ f(x_n) \}_{n=1}^{\infty} \) does not converge to \( f(c) \). Then \( f \) is \hl{discontinuous} at \( c \).}
    \item \hl{\textbf{Theorem 3.3.2 (Minimum-maximum theorem / Extreme value theorem).}} \textit{A continuous function \( f: [a, b] \to \mathbb{R} \) achieves both an absolute minimum and an absolute maximum on \( [a, b] \).}
    \item \hl{\textbf{Theorem 3.3.8 (Bolzano’s intermediate value theorem).}} \textit{Let \( f: [a, b] \to \mathbb{R} \) be a continuous function. Suppose \( y \in \mathbb{R} \) is such that \( f(a) < y < f(b) \) or \( f(a) > y > f(b) \). Then there exists a \( c \in (a, b) \) such that \( f(c) = y \).}
    \item \textbf{Proposition 3.3.10.} \textit{Let \( f(x) \) be a polynomial of odd degree. Then \( f \) has a real root.}
    \item \hl{\textbf{Definition 3.4.1.}} Let \( S \subset \mathbb{R} \), and let \( f: S \to \mathbb{R} \) be a function. Suppose for every \( \epsilon > 0 \) there exists a \( \delta > 0 \) such that whenever \( x, c \in S \) and \( |x - c| < \delta \), then \( |f(x) - f(c)| < \epsilon \). Then we say \( f \) is \hl{\textit{uniformly continuous}}.
    \item \textbf{Theorem 3.4.4.} \textit{Let \( f : [a, b] \to \mathbb{R} \) be a continuous function. Then \( f \) is uniformly continuous.}
    \item \textbf{Lemma 3.4.5.} \textit{Let \( S \subset \mathbb{R} \) and let \( f : S \to \mathbb{R} \) be a uniformly continuous function. Let \( \{x_n\}_{n=1}^{\infty} \) be a Cauchy sequence in \( S \). Then \( \{f(x_n)\}_{n=1}^{\infty} \) is Cauchy.}
    \item \hl{\textbf{Definition 3.4.7.}} \textit{A function \( f : S \to \mathbb{R} \) is \hl{\textit{Lipschitz continuous}}, if there exists a \( K \in \mathbb{R} \), such that}
    \[
    |f(x) - f(y)| \leq K |x - y| \quad \text{for all } x \text{ and } y \text{ in } S.
    \]
    \item \textbf{Proposition 3.4.8.} \textit{A Lipschitz continuous function is uniformly continuous.}
\end{itemize}

\section*{Derivative}
\begin{itemize}
    \item \hl{\textbf{Definition 4.1.1.}} Let \( I \) be an interval, let \( f : I \to \mathbb{R} \) be a function, and let \( c \in I \). If the limit
    \[
    L := \lim_{x \to c} \frac{f(x) - f(c)}{x - c}
    \]
    exists, then we say \( f \) is \textit{differentiable} at \( c \), we call \( L \) the \textit{derivative} of \( f \) at \( c \), and we write
    \[
    f'(c) := L.
    \]
    
    If \( f \) is differentiable at all \( c \in I \), then we simply say that \( f \) is \textit{differentiable}, and then we obtain a function \( f' : I \to \mathbb{R} \). The derivative is sometimes written as \( \frac{df}{dx} \) or \( \frac{d}{dx}(f(x)) \).
    \item Prop: Linearity, product, quotient, chain rule
    \item \hl{\textbf{Definition 4.2.1.}} \textit{Let \( S \subset \mathbb{R} \) be a set and let \( f : S \to \mathbb{R} \) be a function. The function \( f \) is said to have a \hl{\textit{relative maximum}} at \( c \in S \) if there exists a \( \delta > 0 \) such that for all \( x \in S \) where \( |x - c| < \delta \), we have \( f(x) \leq f(c) \). The definition of \textit{relative minimum} is analogous.}
    \item \textbf{Lemma 4.2.2.} \textit{Suppose \( f : (a, b) \to \mathbb{R} \) is differentiable at \( c \in (a, b) \), and \( f \) has a relative minimum or a relative maximum at \( c \). Then \( f'(c) = 0 \).}
    \item \hl{\textbf{Theorem 4.2.3 (Rolle).}} \textit{Let \( f : [a, b] \to \mathbb{R} \) be a continuous function differentiable on \( (a, b) \) such that \( f(a) = f(b) \). Then there exists a \( c \in (a, b) \) such that \( f'(c) = 0 \).}
    \item \hl{\textbf{Theorem 4.2.4 (Mean value theorem).}} \textit{Let \( f : [a, b] \to \mathbb{R} \) be a continuous function differentiable on \( (a, b) \). Then there exists a point \( c \in (a, b) \) such that}
    \[
    f(b) - f(a) = f'(c)(b - a).
    \]
    \item \textbf{Theorem 4.2.5 (Cauchy’s mean value theorem).} \textit{Let \( f : [a, b] \to \mathbb{R} \) and \( \varphi : [a, b] \to \mathbb{R} \) be continuous functions differentiable on \( (a, b) \). Then there exists a point \( c \in (a, b) \) such that}
    \[
    (f(b) - f(a)) \varphi'(c) = f'(c)(\varphi(b) - \varphi(a)).
    \]
    \item \textbf{Proposition 4.2.6.} \textit{Let \( I \) be an interval and let \( f : I \to \mathbb{R} \) be a differentiable function such that \( f'(x) = 0 \) for all \( x \in I \). Then \( f \) is constant.}
    \item \textbf{Proposition 4.2.7.} \textit{Let \( I \) be an interval and let \( f : I \to \mathbb{R} \) be a differentiable function.}
    \begin{itemize}
      \item[(i)] \textit{\( f \) is increasing if and only if \( f'(x) \geq 0 \) for all \( x \in I \).}
      \item[(ii)] \textit{\( f \) is decreasing if and only if \( f'(x) \leq 0 \) for all \( x \in I \).}
    \end{itemize}
    \item \textbf{Proposition 4.2.8.} \textit{Let \( I \) be an interval and let \( f : I \to \mathbb{R} \) be a differentiable function.}
    \begin{itemize}
      \item[(i)] \textit{If \( f'(x) > 0 \) for all \( x \in I \), then \( f \) is strictly increasing.}
      \item[(ii)] \textit{If \( f'(x) < 0 \) for all \( x \in I \), then \( f \) is strictly decreasing.}
    \end{itemize}
    \item \textbf{Proposition 4.2.9.} \textit{Let \( f : (a, b) \to \mathbb{R} \) be continuous. Let \( c \in (a, b) \) and suppose \( f \) is differentiable on \( (a, c) \) and \( (c, b) \).}
    \begin{itemize}
      \item[(i)] \textit{If \( f'(x) \leq 0 \) whenever \( x \in (a, c) \) and \( f'(x) \geq 0 \) whenever \( x \in (c, b) \), then \( f \) has an absolute minimum at \( c \).}
      \item[(ii)] \textit{If \( f'(x) \geq 0 \) whenever \( x \in (a, c) \) and \( f'(x) \leq 0 \) whenever \( x \in (c, b) \), then \( f \) has an absolute maximum at \( c \).}
    \end{itemize}
    \item \textbf{Proposition 4.2.10.}
    \begin{itemize}
      \item[(i)] \textit{Suppose \( f : [a, b) \to \mathbb{R} \) is continuous, differentiable in \( (a, b) \), and \( \lim_{x \to a} f'(x) = L \). Then \( f \) is differentiable at \( a \) and \( f'(a) = L \).}
      \item[(ii)] \textit{Suppose \( f : (a, b] \to \mathbb{R} \) is continuous, differentiable in \( (a, b) \), and \( \lim_{x \to b} f'(x) = L \). Then \( f \) is differentiable at \( b \) and \( f'(b) = L \).}
    \end{itemize}
    \item \hl{\textbf{Theorem 4.2.11 (Darboux).}} \textit{Let \( f : [a, b] \to \mathbb{R} \) be differentiable. Suppose \( y \in \mathbb{R} \) is such that \( f'(a) < y < f'(b) \) or \( f'(a) > y > f'(b) \). Then there exists a \( c \in (a, b) \) such that \( f'(c) = y \).}
    \item \hl{\textbf{Definition 4.3.1.}} \textit{For an \( n \) times differentiable function \( f \) defined near a point \( x_0 \in \mathbb{R} \), define the \( n \)th order \textit{Taylor polynomial} for \( f \) at \( x_0 \) as}
    \[
    P_n^{x_0}(x) := \sum_{k=0}^{n} \frac{f^{(k)}(x_0)}{k!}(x - x_0)^k
    \]
    \item \hl{\textbf{Theorem 4.3.2 (Taylor).}} \textit{
    Suppose \( f : [a, b] \to \mathbb{R} \) is a function with \( n \) continuous derivatives on \( [a, b] \) and such that \( f^{(n+1)} \) exists on \( (a, b) \). Given distinct points \( x_0 \) and \( x \) in \( [a, b] \), we can find a point \( c \) between \( x_0 \) and \( x \) such that
    \[
    f(x) = P_n^{x_0}(x) + \frac{f^{(n+1)}(c)}{(n+1)!}(x - x_0)^{n+1}.
    \]
    }
    \item \textbf{Proposition 4.3.3 (Second derivative test).} \textit{
    Suppose \( f : (a, b) \to \mathbb{R} \) is twice continuously differentiable, \( x_0 \in (a, b) \), \( f'(x_0) = 0 \) and \( f''(x_0) > 0 \). Then \( f \) has a strict relative minimum at \( x_0 \).}
    \item \textbf{Lemma 4.4.1.}
    \textit{Let } $I, J \subset \mathbb{R}$ \textit{ be intervals. If } $f : I \to J$ \textit{ is strictly monotone (hence one-to-one), onto (}$f(I) = J$\textit{), differentiable at } $x_0 \in I$, \textit{and } $f'(x_0) \neq 0$, \textit{then the inverse } $f^{-1}$ \textit{ is differentiable at } $y_0 = f(x_0)$ \textit{ and}
    \[
    (f^{-1})'(y_0) = \frac{1}{f'(f^{-1}(y_0))} = \frac{1}{f'(x_0)}.
    \]    
    \textit{If } $f$ \textit{ is continuously differentiable and } $f'$ \textit{ is never zero, then } $f^{-1}$ \textit{ is continuously differentiable.}
    \item \hl{\textbf{Theorem 4.4.2 (Inverse function theorem).}}
    \textit{
    Let $f : (a, b) \to \mathbb{R}$ be a continuously differentiable function, $x_0 \in (a, b)$ a point where $f'(x_0) \neq 0$. Then there exists an open interval $I \subset (a, b)$ with $x_0 \in I$, the restriction $f|_I$ is injective with a continuously differentiable inverse $g : J \to I$ defined on an interval $J := f(I)$, and
    \[
    g'(y) = \frac{1}{f'(g(y))} \quad \text{for all } y \in J.
    \]
    }
    

\end{itemize}

\section*{Riemann Integral}
\begin{itemize}
    \item \hl{\textbf{Definition 5.1.1.}} A \textit{partition} \( P \) of \( [a, b] \) is a finite set of numbers \( \{ x_0, x_1, x_2, \ldots, x_n \} \) such that
    \[
    a = x_0 < x_1 < x_2 < \cdots < x_{n-1} < x_n = b.
    \]    
    We write
    \[
    \Delta x_i := x_i - x_{i-1}.
    \]    
    Suppose \( f : [a, b] \to \mathbb{R} \) is bounded and \( P \) is a partition of \( [a, b] \). Define
    \[
    m_i := \inf \{ f(x) : x_{i-1} \le x \le x_i \},\]\[ \quad M_i := \sup \{ f(x) : x_{i-1} \le x \le x_i \},
    \]
    \[
    L(P, f) := \sum_{i=1}^n m_i \Delta x_i, \quad U(P, f) := \sum_{i=1}^n M_i \Delta x_i.
    \]  
    We call \( L(P, f) \) the \textit{lower Darboux sum} and \( U(P, f) \) the \textit{upper Darboux sum}.

    \item \textbf{Proposition 5.1.2.} Let \( f : [a, b] \to \mathbb{R} \) be a bounded function. Let \( m, M \in \mathbb{R} \) be such that for all \( x \in [a, b] \), we have \( m \le f(x) \le M \). Then for every partition \( P \) of \( [a, b] \),
    \[
    m(b - a) \le L(P, f) \le U(P, f) \le M(b - a).\]

    \item \hl{\textbf{Definition 5.1.3.}} As the sets of lower and upper Darboux sums are bounded, we define
    \[
    \underline{\int_a^b} f(x) \, dx := \sup \{ L(P, f) : P \text{ a partition of } [a, b] \},
    \]
    \[
    \overline{\int_a^b} f(x) \, dx := \inf \{ U(P, f) : P \text{ a partition of } [a, b] \}.
    \]
    We call \( \underline{\int} \) the \textit{lower Darboux integral} and \( \overline{\int} \) the \textit{upper Darboux integral}. To avoid worrying about the variable of integration, we often simply write
    \[
    \underline{\int_a^b} f := \underline{\int_a^b} f(x) \, dx
    \quad \text{and} \quad
    \overline{\int_a^b} f := \overline{\int_a^b} f(x) \, dx.
    \]

    \item \textbf{Definition 5.1.6.} Let \( P = \{x_0, x_1, \ldots, x_n\} \) and \( \widetilde{P} = \{\widetilde{x}_0, \widetilde{x}_1, \ldots, \widetilde{x}_\ell\} \) be partitions of \([a,b]\). We say \( \widetilde{P} \) is a \textit{refinement} of \( P \) if as sets \( P \subset \widetilde{P} \).

    That is, \( \widetilde{P} \) is a refinement of a partition if it contains all the numbers in \( P \) and perhaps some other numbers in between. For example, \(\{0, 0.5, 1, 2\}\) is a partition of \([0,2]\) and \(\{0, 0.2, 0.5, 1, 1.5, 1.75, 2\}\) is a refinement. The main reason for introducing refinements is the following proposition.
    
    \item \textbf{Proposition 5.1.7.} \textit{Let \( f : [a,b] \to \mathbb{R} \) be a bounded function, and let \( P \) be a partition of \([a,b]\). Let \( \widetilde{P} \) be a refinement of \( P \). Then}
    \[
    L(P,f) \leq L(\widetilde{P}, f) \quad \text{and} \quad U(\widetilde{P}, f) \leq U(P, f).
    \]

    \item \textbf{Proposition 5.1.8.} \textit{Let \( f : [a,b] \to \mathbb{R} \) be a bounded function. Let \( m, M \in \mathbb{R} \) be such that for all \( x \in [a,b] \), we have \( m \leq f(x) \leq M \). Then}
    \[
    m(b - a) \leq \underline{\int_a^b} f \leq \overline{\int_a^b} f \leq M(b - a).
    \]

    \item \textbf{Definition 5.1.9.} Let \( f : [a,b] \to \mathbb{R} \) be a bounded function such that
    \[
    \underline{\int_a^b} f(x)\,dx = \overline{\int_a^b} f(x)\,dx.
    \]
    Then \( f \) is said to be \textit{Riemann integrable}. The set of Riemann integrable functions on \( [a,b] \) is denoted by \( \mathfrak{R}([a,b]) \). When \( f \in \mathfrak{R}([a,b]) \), we define
    \[
    \int_a^b f(x)\,dx := \underline{\int_a^b} f(x)\,dx = \overline{\int_a^b} f(x)\,dx.
    \]
    As before, we often write
    \[
    \int_a^b f := \int_a^b f(x)\,dx.
    \]
    The number \( \int_a^b f \) is called the \textit{Riemann integral} of \( f \), or sometimes simply the \textit{integral} of \( f \).

    \item \textbf{Proposition 5.1.10.} Let \( f : [a,b] \to \mathbb{R} \) be a Riemann integrable function. Let \( m, M \in \mathbb{R} \) be such that \( m \leq f(x) \leq M \) for all \( x \in [a,b] \). Then
    \[
    m(b-a) \leq \int_a^b f \leq M(b-a).
    \]

    \item \hl{\textbf{Proposition 5.1.13.}} \textit{Let } f : [a, b] \( \to \mathbb{R} \) \textit{ be a bounded function. Then } f \textit{ is Riemann integrable if for every } \( \varepsilon > 0 \), \textit{ there exists a partition } P \textit{ of } [a, b] \textit{ such that}
    \[
    U(P, f) - L(P, f) < \varepsilon.
    \]

    \item \hl{\textbf{Lemma 5.2.7.}} If \( f : [a,b] \to \mathbb{R} \) is a continuous function, then \( f \in \mathcal{R}([a,b]) \).

    \item \textbf{Theorem 5.2.9.} Let \( f : [a,b] \to \mathbb{R} \) be a bounded function with finitely many discontinuities. Then \( f \in \mathcal{R}([a,b]) \).

    \item \hl{\textbf{Theorem 5.3.1.}} 
    Let \( F : [a, b] \to \mathbb{R} \) be a continuous function, differentiable on \( (a, b) \). 
    Let \( f \in \mathcal{R}([a, b]) \) be such that \( f(x) = F'(x) \) for \( x \in (a, b) \). Then
    \[
    \int_a^b f = F(b) - F(a).
    \]

    \item \hl{\textbf{Theorem 5.3.3.} } 
    Let \( f : [a, b] \to \mathbb{R} \) be a Riemann integrable function. Define  
    \[
    F(x) := \int_a^x f.
    \]  
    First, \( F \) is continuous on \( [a, b] \).  
    Second, if \( f \) is continuous at \( c \in [a, b] \), then \( F \) is differentiable at \( c \) and \( F'(c) = f(c) \).

    








\end{itemize}
\end{multicols}

\end{document}
